\let\negmedspace\undefined
\let\negthickspace\undefined
\documentclass[journal]{IEEEtran}
\usepackage[a5paper, margin=10mm, onecolumn]{geometry}
%\usepackage{lmodern} 
\usepackage{tfrupee}

\setlength{\headheight}{1cm}
\setlength{\headsep}{0mm}

\usepackage{gvv-book}
\usepackage{gvv}
\usepackage{cite}
\usepackage{amsmath,amssymb,amsfonts,amsthm}
\usepackage{algorithmic}
\usepackage{graphicx}
\usepackage{textcomp}
\usepackage{xcolor}
\usepackage{txfonts}
\usepackage{listings}
\usepackage{enumitem}
\usepackage{mathtools}
\usepackage{gensymb}
\usepackage{comment}
\usepackage[breaklinks=true]{hyperref}
\usepackage{tkz-euclide} 
\usepackage{listings}                                      
\def\inputGnumericTable{}                                 
\usepackage[latin1]{inputenc}                                
\usepackage{color}                                            
\usepackage{array}                                            
\usepackage{longtable}                                       
\usepackage{calc}                                             
\usepackage{multirow}                                         
\usepackage{hhline}                                           
\usepackage{ifthen}                                           
\usepackage{lscape}
\usepackage{circuitikz}
\tikzstyle{block} = [rectangle, draw, fill=blue!20, 
    text width=4em, text centered, rounded corners, minimum height=3em]
\tikzstyle{sum} = [draw, fill=blue!10, circle, minimum size=1cm, node distance=1.5cm]
\tikzstyle{input} = [coordinate]
\tikzstyle{output} = [coordinate]

\begin{document}

\bibliographystyle{IEEEtran}
\vspace{3cm}

\title{4.13.14}
\author{EE25BTECH11009 - Anshu Kumar Ram }
\maketitle
{\let\newpage\relax\maketitle}

\renewcommand{\thefigure}{\theenumi}
\renewcommand{\thetable}{\theenumi}
\setlength{\intextsep}{10pt}

\numberwithin{equation}{enumi}
\numberwithin{figure}{enumi}
\renewcommand{\thetable}{\theenumi}

\textbf{Question}:\\
The perpendicular bisector of the line segment joining $P(1,4)$ and $Q(k,3)$ has $Y$--intercept $-4$. Then a possible value of $k$ is
\begin{align}
\text{(a) }1 \quad \text{(b) }2 \quad \text{(c) }-2 \quad \text{(d) }-4
\end{align}

\solution \\

Equation of perpendicular bisector:
\begin{align}
	\vec{n}^{\top}\vec{x} = c \label{eq:1}
\end{align}
where direction vector of $PQ$ is
\begin{align}
	\vec{PQ} = \myvec{k-1 \\ -1}, \label{eq:2}
\end{align}
As the perpendicular bisector has normal,
\begin{align}
\vec{n} = \myvec{k-1 \\ -1}.
\end{align}
The midpoint of $PQ$ is
\begin{align}
\vec{M} = \myvec{\tfrac{1+k}{2} \\[4pt] \tfrac{7}{2}}.
\end{align}
And 
\begin{align}
	c=-4 \label{eq:3}
\end{align}

From ~\eqref{eq:2} and ~\eqref{eq:3} in ~\eqref{eq:1}, 
\begin{align}
\myvec{k-1 & -1}\myvec{-\tfrac{1+k}{2} \\[4pt] -4-\tfrac{7}{2}}=0
\end{align}

\begin{align}
\therefore -(k^2-1) + 15 = 0
\end{align}

\begin{align}
-k^2+16=0 \quad \Rightarrow \quad k^2=16
\end{align}

Hence
\begin{align}
k = \pm 4
\end{align}

From the given options, the valid value is
\[
\boxed{-4}
\]

\begin{figure}[h!]
    \centering
    \includegraphics[height=0.5\textheight, keepaspectratio]{figs/Figure_1.png}
    \label{figure_1}
\end{figure}

\end{document}
